% Part: formatives
% Chapter: one
% Section: propositional-modal-logic

\documentclass[../../../include/open-logic-section]{subfiles}

\begin{document}

\olfileid{for}{one}{pmod}

\newpage
\olsection{Propositional Modal Logic}

\begin{prob}
 Consider the model $\mModel{M}$ with:
 \begin{itemize}
    \item Worlds $W=\{@,u,v\}$, where $@$ is the actual world,
    \item Accessibility relation 
    $R=\{\tuple{@,u},\tuple{@,v},\tuple{u,@},\tuple{v,u},\tuple{v,v}\}$
    \item Valuation $V(\Obj{\pvar{P}},@)=V(\Obj{\pvar{Q}},u)=V(\Obj{\pvar{R}},v)=\True$, all else $\False$.
 \end{itemize}
 Say whether the formulas below are true in the model (i.e., true at
    $@$). Briefly justify your answers.
    \begin{enumerate}
    \item $\Diamond \Obj{\pvar{P}}$
    \item $\Box(\Obj{\pvar{R}}\lif\Diamond \Obj{\pvar{Q}})$
    \item $\Box(\Diamond \Obj{\pvar{R}}\lor\Diamond \Obj{\pvar{P}})$
    \item $\Box\Box\Diamond \Obj{\pvar{Q}}$
    \end{enumerate}

\begin{ans}
\item $\Diamond \Obj{\pvar{P}}$. False. Justification: $@$ only has
access to worlds where $\Obj{\pvar{P}}$ is false.
\item $\Box(\Obj{\pvar{R}}\lif\Diamond \Obj{\pvar{Q}})$. True.
Justification: $\Obj{\pvar{R}}$ is false at $u$, $\Diamond
\Obj{\pvar{Q}}$ is true at $v$ (since $v$ has access to world $u$
where $\Obj{\pvar{Q}}$ is true), therefore $\Obj{\pvar{R}}\lif\Diamond
\Obj{\pvar{Q}}$ is true at $u$ and $v$. Since these are the only two
worlds $@$ has access to, $\Box(\Obj{\pvar{R}}\lif\Diamond
\Obj{\pvar{Q}})$ is true at $@$, hence true.
\item $\Box(\Diamond \Obj{\pvar{R}}\lor\Diamond \Obj{\pvar{P}})$. True. 
Justification: $\Diamond\Obj{\pvar{P}}$ is true at $u$ (since $u$ has access
to world $@$ where $\Obj{\pvar{P}}$ is true) and 
$\Diamond\Obj{\pvar{R}}$ is true at $v$ (since $v$ has access to 
itself where $\Obj{\pvar{R}}$ is true). Hence 
$\Diamond\Obj{\pvar{R}}\lor\Diamond\Obj{\pvar{P}}$ is true at both $u,v$.
Since these are the only two worlds $@$ has access to, 
$\Box(\Diamond\Obj{\pvar{R}}\lor\Diamond\Obj{\pvar{P}})$ holds at $@$, 
hence is true.
\item $\Box\Box\Diamond \Obj{\pvar{Q}}$. False. Justification:
$\Diamond\Obj{\pvar{Q}}$ is false at $u$, because $u$ doesn't 
have access to a world where $\Obj{\pvar{Q}}$ is true (even though 
$\Obj{\pvar{Q}}$ is true at $u$, $u$ doesn't have access to $u$, 
it only has access to $@$ where $\Obj{\pvar{Q}}$ isn't true). Since
$\Diamond\Obj{\pvar{Q}}$ is false at $u$ and $v$ has access to $u$, 
$\Box\Diamond\Obj{\pvar{Q}}$ is false at $v$. Since 
$\Box\Diamond\Obj{\pvar{Q}}$ is false at $v$ and $@$ has access to 
$v$, $\Box\Box\Diamond\Obj{\pvar{Q}}$ is false at $@$.

\clearpage
\end{ans}

\end{prob}

\begin{prob}
    Establish the following. Use countermodels for invalidities, 
    natural deduction !!{derivation}s for validities. 
    \begin{enumerate}
        \item $\Box\lnot\Obj{\pvar{A}}, 
        \Box(\Obj{\pvar{A}}\lor\Obj{\pvar{B}})
        \Entails[\Log{K}] \Box{\Obj{\pvar{B}}}$

        \item $\Diamond\lnot\Obj{\pvar{A}}, 
        \Box(\Obj{\pvar{A}}\lor\Obj{\pvar{B}})
        \Entails/[\Log{T}] \Box{\Obj{\pvar{B}}}$

        \item $\Entails[\Log{T}]\Diamond(
            (\Obj{\pvar{A}}\land\Obj{\pvar{B}})\lif\Obj{\pvar{A}})$

        \item $\Diamond\Diamond \Obj{\pvar{A}},
        \Box(\Obj{\pvar{B}}\lif\lnot\Obj{\pvar{A}})
        \Entails[\Log{S4}] \lnot\Box\Obj{\pvar{B}}$
    
    \end{enumerate}

    \begin{ans}
    \begin{enumerate}
        \item Proof of $\Box\lnot\Obj{\pvar{A}}, 
        \Box(\Obj{\pvar{A}}\lor\Obj{\pvar{B}})
        \Proves[\Log{K}] \Box{\Obj{\pvar{B}}}$

        \begin{prooftree}
        \AxiomC{$\Box\lnot\Obj{\pvar{A}}$}
            \AxiomC{$\Box(\Obj{\pvar{A}}\lor\Obj{\pvar{B}})$}
                \AxiomC{$\Discharge{\lnot\Obj{\pvar{A}}}{1}$}
                    \AxiomC{$\Discharge{\Obj{\pvar{A}}\lor\Obj{\pvar{B}}}{1}$}
                \RightLabel{\Log{PL}}
                \BinaryInfC{$\Obj{\pvar{B}}$}
        \DischargeRule{$\Box$K}{1}
        \TrinaryInfC{$\Box\Obj{\pvar{B}}$}
        \end{prooftree}

        \emph{Pedantic note.} I write ``proof of
        $\ldots\Proves\ldots$'' rather than ``proof of
        $\ldots\Entails\ldots$'' since a proof primarily establishes
        that the argument is \emph{provable} ($\Proves$) rather than
        \emph{valid} ($\Entails$). However, because our proof systems are
        sounds for their respective semantics, if an argument is provable 
        in a modal system it is valid in that modal system.

        \item Counterexample to establish $\Diamond\lnot\Obj{\pvar{A}}, 
        \Box(\Obj{\pvar{A}}\lor\Obj{\pvar{B}})
        \Entails/[\Log{T}] \Box{\Obj{\pvar{B}}}$

        Model $\mModel{M}$ with $W=\{@,u\}$,
        $R=\{\tuple{@,@},\tuple{@,u},\tuple{u,u}\}$,
        valuation $V(\Obj{\pvar{A}},@)=V(\Obj{\pvar{B}},u)=\True$,
        $V(\Obj{\pvar{A}},u)=V(\Obj{\pvar{B}},@)=\False$, all else
        $\False$.

        Since $\Obj{\pvar{A}}$ is true at $@$ and $\Obj{\pvar{B}}$ is
        true at $v$, $\Obj{\pvar{A}\lor\pvar{B}}$ is true at $@,u$.
        Since these are the only worlds $@$ has access to,
        $\Box(\Obj{\pvar{A}}\lor\Obj{\pvar{B}})$ is true at $@$. Since
        $\Obj{\pvar{A}}$ is false at $u$, $\Obj{\lnot\pvar{A}}$ is
        true at $u$; since $@$ has access to $u$,
        $\Diamond\Obj{\lnot\pvar{A}}$ is true at $@$. However, since
        $\Obj{\pvar{B}}$ is not true at $@$ and $@$ has access to $@$,
        $\Box\Obj{\pvar{B}}$ is not true at $@$.
        
        \emph{Comment. This is the simplest kind of counterexample. It
        would also work if you made $\Obj{\pvar{A}}$ false at $@$ and
        true at $u$ and $\Obj{\pvar{B}}$ true at $@$ and false at $u$
        instead. What's required is to have at least one of
        $\Obj{\pvar{A}},\Obj{\pvar{B}}$ true at every world accessible
        from $@$ (to make $\Box(\Obj{\pvar{A}}\lor\Obj{\pvar{B}})$
        true), but $\Obj{\pvar{A}}$ false in one of these worlds (to
        make $\Diamond\lnot\Obj{\pvar{A}}$ true) and $\Obj{\pvar{B}}$
        false at another one of these worlds (to make
        $\Box\Obj{\pvar{B}}$ false).}

        
        \item Proof of $\vdash_{\Log{T}}\Diamond(
            (\Obj{\pvar{A}}\land\Obj{\pvar{B}})\lif\Obj{\pvar{A}})$

        \begin{prooftree}
        \AxiomC{}
        \RightLabel{\Log{PL}}
        \UnaryInfC{$(\Obj{\pvar{A}}\land\Obj{\pvar{B}})\lif\Obj{\pvar{A}}$}
        \RightLabel{$\Diamond$T}
        \UnaryInfC{$\Diamond((\Obj{\pvar{A}}\land\Obj{\pvar{B}})\lif\Obj{\pvar{A}})$}
        \end{prooftree}

        \emph{Comment. Note that the proof has no premise, since we 
        can derive $(\Obj{\pvar{A}}\land\Obj{\pvar{B}})\lif\Obj{\pvar{A}}$
        in propositional logic without premise, as follows:}
        \begin{prooftree}
        \AxiomC{$\Discharge{\Obj{\pvar{A}}\land\Obj{\pvar{B}}}{1}$}
        \RightLabel{\Elim{\land}}
        \UnaryInfC{$\Obj{\pvar{A}}$}
        \DischargeRule{\Intro{\lif}}{1}
        \UnaryInfC{$(\Obj{\pvar{A}}\land\Obj{\pvar{B}})\lif\Obj{\pvar{A}}$}
        \end{prooftree}
        \emph{We needn't spell that out. We simply use a \Log{PL} 
        step. You don't even need to haved worked out the proof 
        above: it's enough for you to realize that the formula is a 
        logical truth of propositional logic; you then know 
        (by soundness) that there will be a way to prove it.}

        If you only know the $\Box$T rule, you can use this proof
        instead:
        \begin{prooftree}
        \AxiomC{$\Discharge{\Box\lnot((\Obj{\pvar{A}}\land\Obj{\pvar{B}})\lif\Obj{\pvar{A}})}{2}$}
            \AxiomC{$\Discharge{\lnot((\Obj{\pvar{A}}\land\Obj{\pvar{B}})\lif\Obj{\pvar{A}})}{1}$}
            \RightLabel{\Log{PL}}
            \UnaryInfC{$\bot$}
        \DischargeRule{$\Box$K}{1}
        \BinaryInfC{$\Box\bot$}
        \RightLabel{$\Box$T}
        \UnaryInfC{$\bot$}
        \DischargeRule{\Log{PL} ($\bot_C$)}{2}
        \UnaryInfC{$\lnot\Box\lnot((\Obj{\pvar{A}}\land\Obj{\pvar{B}})\lif\Obj{\pvar{A}})$}
        \RightLabel{Dual}
        \UnaryInfC{$\Diamond((\Obj{\pvar{A}}\land\Obj{\pvar{B}})\lif\Obj{\pvar{A}})$}
        \end{prooftree}

        \emph{Comment. The strategy is to establish the conclusion
        $\Diamond!A$ by establishing $\lnot\Box\lnot!A$. We aim to
        establish the later by \emph{reductio}: assume $\Box\lnot!A$
        and derive a contradiction. We note that our $!A$ (that is,
        $(\Obj{\pvar{A}}\land\Obj{\pvar{B}})\lif\Obj{\pvar{A}}$ is a
        tautology), hence $\lnot!A$ a contradiction. Therefore we know
        that $\lnot!A$ we can derive $\bot$: we needn't work out a
        proof, but we can assume that there is one and use a
        $\Log{PL}$ step instead. Using that fact, we get $\Box\bot$
        from $\Box\lnot!A$. That is not yet a contradiction (it
        wouldn't be in the \Log{K} system), but in the \Log{T} system
        we can then derive $\bot$ itself, completing our
        \emph{reductio}. Note that the $\bot_C$ step isn't merely
        deriving anything from $\bot$ (the Ex Falso Quodlibet rule),
        but discharging an assumption and inferring its negation (the
        negation elimination or Reductio rule).}
        
        \item Proof of $\Diamond\Diamond \Obj{\pvar{A}},
        \Box(\Obj{\pvar{B}}\lif\lnot\Obj{\pvar{A}})
        \Proves[\Log{S4}] \lnot\Box\Obj{\pvar{B}}$

        One strategy, using $\Diamond$K to prove
        $\Diamond\lnot\Obj{\pvar{B}}$ first, then the conclusion by
        Duality:

        \begin{prooftree}
            \AxiomC{$\Diamond\Diamond\Obj{\pvar{A}}$}
            \RightLabel{$\Diamond$4}
            \UnaryInfC{$\Diamond\Obj{\pvar{A}}$}
                \AxiomC{$\Box(\Obj{\pvar{B\lif\lnot\pvar{A}}})$}
                    \AxiomC{$\Discharge{\Obj{\pvar{A}}}{1}$}
                        \AxiomC{$\Discharge{\Obj{\pvar{B\lif\lnot\pvar{A}}}}{1}$}
                        \RightLabel{\Log{PL} (\Elim{\lif})}
                    \BinaryInfC{$\lnot\Obj{\pvar{B}}$}
            \DischargeRule{$\Diamond$K}{1}
            \TrinaryInfC{$\Diamond\lnot\Obj{\pvar{B}}$}
            \RightLabel{Dual}
            \UnaryInfC{$\lnot\Box\Obj{\pvar{B}}$}
            \end{prooftree}
    
        Another strategy, \emph{reductio} of $\Box\Obj{\pvar{B}}$:
        \begin{prooftree}
            \AxiomC{$\Discharge{\Box\Obj{\pvar{B}}}{2}$}
            \AxiomC{$\Box(\Obj{\pvar{B}}\lif\lnot\Obj{\pvar{A}})$}
                \AxiomC{$\Discharge{\Obj{\pvar{B}}}{1}$}
                    \AxiomC{$\Discharge{\Obj{\pvar{B}}\lif\lnot\Obj{\pvar{A}}}{1}$}
                \RightLabel{\Log{PL}}
                \BinaryInfC{$\lnot\Obj{\pvar{A}}$}
            \DischargeRule{$\Box$K}{1}
            \TrinaryInfC{$\Box\lnot\Obj{\pvar{A}}$}
            \RightLabel{Dual}
            \UnaryInfC{$\lnot\Diamond\Obj{\pvar{A}}$}
                \AxiomC{$\Diamond\Diamond\Obj{\pvar{A}}$}
                \RightLabel{$\Diamond$4}
                \UnaryInfC{$\Diamond\Obj{\pvar{A}}$}
            \DischargeRule{\Log{PL} ($\bot_C$)}{2}
            \BinaryInfC{$\lnot\Box\Obj{\pvar{B}}$}
        \end{prooftree}
    
        If you do nor know the $\Diamond$4 or $\Diamond$K rules,
        follow the \emph{reductio} strategy above but derive
        $\Box\Box\lnot\Obj{\pvar{A}}$ first, then transform it into
        $\lnot\Diamond\Diamond\Box\Obj{\pvar{B}}$ using the $\Box$K
        and Dual rules like so:
        {\footnotesize
        \begin{prooftree}
            \AxiomC{$\Discharge{\Box\Obj{\pvar{B}}}{3}$}
            \AxiomC{$\Box(\Obj{\pvar{B}}\lif\lnot\Obj{\pvar{A}})$}
                \AxiomC{$\Discharge{\Obj{\pvar{B}}}{1}$}
                    \AxiomC{$\Discharge{\Obj{\pvar{B}}\lif\lnot\Obj{\pvar{A}}}{1}$}
                \RightLabel{\Log{PL}}
                \BinaryInfC{$\lnot\Obj{\pvar{A}}$}
            \DischargeRule{$\Box$K}{1}
            \TrinaryInfC{$\Box\lnot\Obj{\pvar{A}}$}
            \RightLabel{$\Box$4}
            \UnaryInfC{$\Box\Box\lnot\Obj{\pvar{A}}$}
                \AxiomC{$\Discharge{\Box\lnot\Obj{\pvar{A}}}{2}$}
                \RightLabel{Dual}
                \UnaryInfC{$\lnot\Diamond\Obj{\pvar{A}}$}
                \DischargeRule{$\Box$K}{2}
            \BinaryInfC{$\Box\lnot\Diamond\Obj{\pvar{A}}$}
            \RightLabel{Dual}
            \UnaryInfC{$\lnot\Diamond\Diamond\Obj{\pvar{A}}$}
                \AxiomC{$\Diamond\Diamond\Obj{\pvar{A}}$}
            \DischargeRule{\Log{PL}}{3}
            \BinaryInfC{$\lnot\Box\Obj{\pvar{B}}$}
        \end{prooftree}
        }
     
    \end{enumerate}

    \clearpage
    \end{ans}

\end{prob}

\begin{prob}
    We say that relation $R$ has \emph{secondary reflexivity} iff 
    whenever something has that relation to something, the latter has
    that relation to itself: if $Rww'$ then $Rw'w'$. A modal !!{structure}
    is secondary reflexive iff its accessibility relation is secondary 
    reflexive.

    \begin{enumerate}
    \item Can a model can be secondary reflexive without 
    being serial? Briefly justify your answer.
    \item Show that any transitive and symmetric model is secondary 
    reflexive, but not the other way round.
    \item Can you give a characteristic schema for secondary reflexive
    models? Show that the schema is valid for any secondary reflexive
    models. Optional: show that it is not valid on frames that are not
    secondary reflexive.
    \end{enumerate}

    \begin{ans}
    \begin{enumerate}
    \item Yes, a model can be secondary reflexive without 
    being serial. All secondary reflexivity requires is that 
    \emph{if} a world has access to some world, then the latter
    has access to itself. Hence, for instance, a model where 
    no world has $R$ to any world is vacuously secondary 
    reflexive but not serial. 

    \item Let $\mModel{M}$ be symmetric and transitive and 
    suppose $Rwu$ for some $w,u$. Since $R$ is symmetric, 
    we have $Ruw$. Since $Ruw$, $Rwu$ and $R$ is transitive,
    we have $Ruu$. Generalizing over $w,u$, $R$ is 
    secondary reflexive.

    \emph{Comment. ``Generalizing over $w,u$'' is a way of 
    saying that, because $w,u$ were arbitrary picked, the 
    above holds for any $w,u$. Thus, we have shown that 
    for any $w,u$, if $w$ has access to $u$, then $u$ 
    has access to $u$, which is the definition of 
    secondary reflexivity.}

    \item Schema $\Box(\Box!A\lif!A)$ or $\Box\Box!A\lif \Box!A$
    is characteristic of secondary reflexivity. 
    
    \emph{Comment. Both are correct; they are in fact equivalent 
    in \Log{K} and stronger systems. The idea is this: call a world 
    `factive' iff whatever is necessary at that world is true at that 
    world. With reflexivity, every world is factive. With secondary
    reflexivity, not every world need be factive (so $\Box!A\lif!A$ may not
    be true for every $!A$ at every world) but it must be that 
    every world \emph{that is accessed} is factive. Hence any world 
    we see (from the actual world, or some given world) is a world
    where $\Box!A\lif!A$ holds. Hence the characteristic formulas above.}
    
    \emph{Some \emph{incorrect} answers: $\Box!A\lif!A$,
    $\Diamond!A\lif\Box\Diamond!A$, $\Box\Diamond!A\lif\Diamond!A$.}

    Show that $\Box(\Box!A\lif!A)$ holds in secondary reflexive
    models. Suppose $\mModel{M}$ is secondary reflexive, and let $!A$ be
    any formula and $w$ any world. Suppose for reductio that
    $\Box(\Box!A\lif!A)$ is false at $w$. Then there is some world
    $wu$ such that $Rwu$ and $\Box!A\lif!A$ is false at $u$. $\Box!A$
    is true at $u$ and $!A$ is not true at $u$: hence $u$ doesn't have
    access to $u$. But by secondary reflexivity, since $Rwu$ we have
    $Ruu$. Contradiction. 

    Alternatively, show that  $\Box\Box!A\lif \Box!A$ holds in secondary reflexive
    models. Suppose $\mModel{M}$ is secondary reflexive, and let $!A$
    be a formula and $!A$ a $w$ any world. Suppose $\Box\Box!A$ is true
    at $w$. Let $u$ be any world such that $Rwu$. Since  $\Box\Box!A$ is true
    at $w$ and we have $Rwu$, $\Box!A$ is true at $u$. Since $Rwu$,
    by secondary reflexivity $Ruu$. Since $\Box!A$ is true at $u$
    and $Ruu$, $!A$ is true at $u$. Generalizing over $u$, 
    $!A$ is true at any world that $w$ has access to. Hence
    $\Box!A$ is true at $w$. Generalizing over $w$, any world 
    is such that if $\Box\Box!A$ is true in it then $\Box!A$ 
    is true too, hence $\Box\Box!A\lif \Box!A$ is true at 
    every world.

    Optional: show that the schema fails in non-secondary
    reflexive \emph{frames}. Let $\mModel{F}$ be a frame 
    that is not secondary reflexive: for some $w,u$ in 
    that frame, we have $Rwu$ but not $Ruu$. Consider the model 
    $\mModel{M}$ based on $\mModel{F}$ such that $\Obj{\pvar{P}}$
    is true at every world expect at $u$. At $u$ we have 
    $\Box\Obj{\pvar{P}}$, since any world $u$ has access to 
    is other than $u$ (for we don't have $Ruu$) and $\Obj{\pvar{P}}$
    is true at all worlds other than $u$. However, $\Obj{\pvar{P}}$ 
    is false at $u$, hence $\Box!A\lif!A$ is false at $u$. 
    Hence the schema $\Box(\Box!A\lif!A)$ fails. 
    
    The proof for schema $\Box\Box!A\lif \Box!A$ is similar.

    \end{enumerate}

    \clearpage
    \end{ans}
    
\end{prob}

\end{document}